\documentclass[a4paper]{article}
\usepackage{titlesec}
\usepackage{geometry}
\geometry{hmargin={2cm,2cm}, vmargin={3.5cm,3.5cm}}
\usepackage{graphicx}
\usepackage[table]{xcolor}
\usepackage{hyperref}
\hypersetup{pdfborder = {0 0 0},
colorlinks,
citecolor=black,
filecolor=black,
linkcolor=black,
urlcolor=blue}
\definecolor{logo}{RGB}{20,99,163}
\usepackage[T1]{fontenc}
\usepackage[utf8]{inputenc}
\usepackage[english,italian]{babel}
\usepackage{fancyhdr}
\usepackage{lastpage}

% creazione comando subsubsubsection
\titleclass{\subsubsubsection}{straight}[\subsection]

\newcounter{subsubsubsection}[subsubsection]
\renewcommand\thesubsubsubsection{\thesubsubsection.\arabic{subsubsubsection}}
\titleformat{\subsubsubsection}
  {\normalfont\normalsize\bfseries}{\thesubsubsubsection}{1em}{}
    \titlespacing*{\subsubsubsection}
    {0pt}{3.25ex plus 1ex minus .2ex}{1.5ex plus .2ex}

\makeatletter
\def\toclevel@subsubsubsection{4}
\def\l@subsubsubsection{\@dottedtocline{4}{7em}{4em}}
\makeatother

\setcounter{secnumdepth}{4}
\setcounter{tocdepth}{4}
%fine creazione comando


%Comandi di impaginazione uguale per tutti i documenti
\pagestyle{fancy}
\lhead{\includegraphics[scale=0.2]{../../../../Images/logo.png}}
%Titolo del documento
\rhead{\large *Verbale Esterno 2021-03-25}
%\rfoot{\thepage}
\cfoot{Pagina \thepage{} di \pageref{LastPage}}
\setlength{\headheight}{35pt}
\renewcommand{\footrulewidth}{0.4pt}


\begin{document}
\begin{titlepage}
    \begin{center}
        \includegraphics{../../../../Images/logo.png}\\
        \vspace{20px}
        \textcolor{logo}{\hrulefill}\\
        \vspace{20px}
        \textbf{\huge\textcolor{logo}{Verbale Esterno 2021-03-25}}\\
        \vspace{10px}
        \textcolor{logo}{\hrulefill}\\
        \vspace{40px}
        \textbf{\Large Informazioni sul documento}\\
        \vspace{20px}
        \begin{tabular}{p{100px} | p{100px}}
            \textbf{Versione} & 0.0.1\\
            \textbf{Responsabile} & Samuele De Simone\\
            \textbf{Redattori} & Giaccone Nicolò\\
            \textbf{Verificatore} & Amedeo Meggiolaro\\
            \textbf{Destinatari} & Prof. Tullio Vardanega \newline Prof. Riccardo Cardin, \newline TechSWEave\\
            \textbf{Stato} & Da approvare\\
            \textbf{Uso} & Esterno\\
        \end{tabular}\\
        \vspace{60px}
        \href{mailto:techsweave@gmail.com}{techsweave@gmail.com}\\

    \end{center}
    \end{titlepage}

    %DIARIO DELLE MODIFICHE
    \begin{center}
        \textbf{\Large Diario delle modifiche}\\
        \vspace{10px}
        \begin{table}[h!]
        \centering
        \rowcolors{2}{logo!10}{logo!40}
        \renewcommand{\arraystretch}{1.8}
        \begin{tabular}{p{150px} p{90px} p{50px} p{60px} p{45px}}
            \rowcolor{logo!70} \textbf{Modifica} & \textbf{Autore} & \textbf{Ruolo} & \textbf{Data} & \textbf{Versione}\\
            Prima stesura del verbale & Giaccone Nicolò & Analista & 2021-03-25 & 0.0.1 \\    
            Verifica del verbale & Amedeo Meggiolaro & Verificatore & 2021-03-26 & 0.1.0 \\    
            Approvazione del verbale & Samuele De Simone & Responsabile & 2021-03-27 & 1.0.0 \\    
        \end{tabular}
    \end{table}
    \end{center}

    %INDICE
    \newpage
    \tableofcontents
    \newpage
    \section{Informazioni riunione}
    \begin{itemize}
        \item Data e ora: 2021/03/25 ore 17:30-18:00.
        \item Luogo: Google Meet fornito dai proponenti.
        \item Partecipanti interni: Amedeo Meggiolaro, Nicolò Giaccone, Manuel Varo, Marco Barbaresco, Samuele De Simone, Simone Urbani, Tito Scutari.
        \item Partecipanti esterni: Alessandro Maccagnan, Milo Ertola.
        \item Segretario: Amedeo Meggiolaro.
    \end{itemize}
    \section{Argomenti della riunione}
    \begin{itemize}
        \item Prosta di alcune funzionalità aggiuntive;
        \item Chiarimento su alcuni requisiti.
    \end{itemize}
    \section{Resoconto}
    
        \subsection{Prosta di alcune funzionalità aggiuntive}
        Sono state  ai proponenti delle funzionalità aggiuntive per la piattaforma. I proponenti ci hanno indicato la fattibilità di queste proposte consigliandoci delle eventuali alternative.
        Le proposte sono le seguenti:
        \begin{itemize}
            \item Servizio di reso: è stato proposto un servizio di reso, i proponenti hanno evidenziato le difficoltà insite in questa proposta. Ci hanno poi indicato delle linee guida da seguire e delle funzionalità obbligatorie che si devono implementare se si vuole sviluppare questa funzionalità;
            \item Recensioni: Il gruppo ha proposto un sistema di recensioni dei prodotti;
            \item Comunicazione tra clienti: Il gruppo ha proposto un servizio di comunicazione tra clienti e venditore, sono state proposte inoltre due modalità di cominicazione
            \begin{itemize}
                \item Esterno: ossia tramite l'invio di mail. Con la presenza di un apposito form e l'utilizzo delle mail associate all'account;
                \item Interno: ossia un servizio di messagistica interna;  
            \end{itemize}
            I proponenti hanno consigliato l'approccio esterno, indicandoci come opzionale l'approccio interno;
        \end{itemize}
        
        \subsection{Chiarimento su alcuni requisiti}
        Il gruppo ha espresso dei dubbi sorti durante l'analisi del capitolato per redarre l'analisi dei requisiti, i proponenti hanno chiarito questi dubbi ed indicato delle linee guida da seguire.
        I dubbi erano i seguenti:
        \begin{itemize}
            \item Servizio per il login: il gruppo ha chiesto chiarimenti su che servizio usare per il login nella piattoforma tra:
            \begin{itemize}
                \item Amazon Cognito: implementazione più semplice, ma poco diffuso e trasparente all'integrazione;
                \item Auth0: implmentazione complicata essendo un servizio esterno ad AWS, molto diffuso.
            \end{itemize}
            \item Strutturazione filtri: il gruppo ha chiesto che filtri di ricerca doveno essere inclusi;
            \item Categorizzazione: il gruppo ha chiesto se le categorie dovessero essere create dal venditore o se dovessero essere pre-esistetni nella piattforma
        \end{itemize}
\end{document}
